\documentclass[11pt,letterpaper]{article}
\usepackage[margin=0.92in,top=0.88in,bottom=0.88in,headheight=14pt]{geometry}
\usepackage{amsmath,amssymb,amsthm,mathtools,booktabs,array,xcolor,fancyhdr,hyperref,microtype}
\usepackage[T1]{fontenc}\usepackage{lmodern}
\definecolor{inkblue}{HTML}{183F54}
\definecolor{accent}{HTML}{176566}
\hypersetup{colorlinks=true,urlcolor=accent,linkcolor=inkblue,citecolor=accent,pdftitle={Sharp local bounds from methods in the Lean proof of Dirichlet GRH},pdfauthor={Research analysis}}
\pagestyle{fancy}\fancyhf{}\lhead{\small\textcolor{inkblue}{Lubin--Tate interpolation and graph towers}}\rhead{\small Research report}\cfoot{\thepage}
\renewcommand{\headrulewidth}{.3pt}
\setlength{\parindent}{1.15em}\setlength{\parskip}{2pt}
\setlength{\emergencystretch}{2.5em}
\newtheorem{theorem}{Theorem}[section]\newtheorem{lemma}[theorem]{Lemma}
\newtheorem{proposition}[theorem]{Proposition}\newtheorem{corollary}[theorem]{Corollary}
\theoremstyle{definition}\newtheorem{example}[theorem]{Example}
\newcommand{\OO}{\mathcal O}\newcommand{\FF}{\mathbb F}\newcommand{\ZZ}{\mathbb Z}\newcommand{\QQ}{\mathbb Q}
\newcommand{\ord}{\operatorname{ord}}\newcommand{\tr}{\operatorname{tr}}
\newcommand{\src}[2]{\href{https://github.com/lsernau/grh-lean/blob/72fead102f663d5c7a407964a13dc0cf9d977207/grh/src/#1}{#2}}
\newcommand{\citesource}[1]{\textcolor{accent}{#1}}
\raggedbottom
\numberwithin{equation}{section}
\begin{document}
\thispagestyle{plain}
{\small\bfseries\textcolor{accent}{NUMBER THEORY \enspace / \enspace RESEARCH REPORT}\hfill 26 November 2026\par}
\vspace{10pt}
{\LARGE\bfseries\color{inkblue} Sharp local bounds from methods in the Lean proof of Dirichlet GRH\par}
\vspace{4pt}
{\large Lubin--Tate interpolation and spanning trees in abelian covers\par}
\vspace{9pt}\hrule\vspace{10pt}
\noindent Two questions in the ProveIt incoming collection admit exact answers using local ideas in the Lean proof. For Lubin--Tate branches, we determine precisely when a proposed precision bound is attained over every $\OO/\pi^r$. We give the optimal common cutoff over all choices of branches and an explicit finite-precision multiplication rule in Berger and Sprang's setting. For rank-two graph towers we bound the growth coefficient for every step count, characterize attainment for even step counts, and compute the unrestricted maximum for four steps at every prime $p>8$. The latter classification singles out seven exceptional primes. Both proofs are independent of the conclusion of Dirichlet GRH.

\section{From the Lean proof to the two questions}
The readable project at commit \texttt{72fead102f66} proves Mathlib's Riemann hypothesis and the corresponding assertions for Dirichlet $L$-functions. Its half-plane theorem quantifies over every character of every positive modulus, excludes the principal pole at $1$, and proves nonvanishing for $\Re s>1/2$. Reflection gives the open-strip statement. The five public declarations can be read in \src{Main.lean}{\texttt{Main.lean}} with their Mathlib predicates in \src{Statements.lean}{\texttt{Statements.lean}}. Andrew Sutherland reports successful independent paranoid comparator runs on both editions, with the usual axioms \texttt{propext}, \texttt{Classical.choice}, and \texttt{Quot.sound} \cite{repo,sutherland}. We inspected pinned sources and used those published verification results.

\paragraph{Why the global argument is specialized.}
The short-vector and near-spanning theorems supply an upper bound for an index-corrected Gram covolume. The chain of theta lifts supplies a lower bound on the same quantity. Cancelling their common main term in
\src{New/Reduction/MainInequality.lean\#L87}{\texttt{New.star\_inequality}} gives
\[
       m\,\mathrm{totalRemainder}\ \le\ f a(\log b)^5.
\]
Here $m$ is the number of lifts; $f$ counts infinite places, $a$ is a derivative order and $b$ is an auxiliary scale. With $m$ between $\sqrt{x}$ and $2\sqrt{x}$, and a prime number theorem at a larger scale, this bounds a smoothed prime sum with weight $\Lambda(n)(1-(n/d))$ on $n\equiv1\bmod q_0$, where $18480\mid q_0$. For every putative zero to the right of the line, the proof chooses an auxiliary quadratic twist free of zeros in a suitable rectangle. The residues in that rectangle then involve sums of nonnegative orders of zeros. A one-sided Landau argument excludes the putative zero. The exact implications are in \src{New/Reduction/OneSidedBound.lean}{\texttt{OneSidedBound.lean}} and \src{New/Reduction/Zeros/Criterion.lean\#L37}{\texttt{Zeros/Criterion.lean}}.

\paragraph{Two local arguments survive a change of problem.}
In \src{New/Derivatives/Sharpness/OneCurve.lean\#L45}{\texttt{exists\_monomial\_max\_norm\_coeff}}, only finitely many monomials can contribute to a specified Taylor coefficient. The ultrametric inequality bounds that coefficient, and a single monomial attains the maximum. For interpolation, we identify all terms that can contribute to the extreme Laurent coefficient; their multiplicity decides equality. In \src{Lib/PadicHodge/TateSen.lean\#L108}{\texttt{exists\_norm\_normalizedTrace\_le}} and \src{Lib/PadicHodge/PeriodRings/AxTheorem.lean\#L347}{\texttt{norm\_zeta\_sub\_one\_pow}}, the valuation of $\zeta-1$ makes a least-order term visible at a cyclotomic point. For the graphs, the same valuation reads the first nonzero directional moment from a Laplacian eigenvalue. The interpolation and cyclotomic principles have independent histories. The archive questions \cite{oldbranches,oldgraphs} specify the settings in which we sharpen them.

\section{How much of each formal branch must one observe?}
Let $K$ be a finite extension of $\QQ_p$, $\OO$ its integers, $\pi$ a uniformizer, $v(\pi)=1$, and $q$ the residue-field size. Fix a Lubin--Tate formal $\OO$-module for $\pi$. Write $[b](X)$ for multiplication by $b\in\OO$. Choose $d\ge1$ distinct directions $B\subset\OO$. At precision $r\ge1$ work over $A_r=\OO/\pi^r$, and put $u_b=[b](X)\bmod\pi^r$. The observation at $b$ consists of the coefficients of $H(X,u_b)$ at degrees $0,\ldots,n_b-1$. The desired output is the list of all coefficients of $H\in A_r[[X,Y]]$ of total degree $<d$.

Berger and Sprang proved existence of a finite common cutoff over the residue field \cite[Lemma 4.2]{bs}. The incoming manuscript \cite{oldbranches} gave a sharp residue-field formula, an implicit exact criterion over $A_r$, and an upper bound there. Its third further question asks which configurations attain the latter bound. We solve that question and determine the optimal common cutoff at every $r$.

\subsection{The extremal coefficient}
For a fixed branch $b$, write $h_c=v(b-c)$ for $c\ne b$. Put
\[
 \omega=\sum_{c\ne b}q^{h_c},\qquad
 m_j=\#\{c\ne b:h_c\ge j\},\qquad H_*=\max_{c\ne b}h_c.
\]
Use $H_*=0$ if the list is empty. For $H_*\ge1$ write $D=(q-1)q^{H_*-1}$ and $M=m_{H_*}$. Let $t_b(r)$ be the least nonnegative $t$ such that $X^t/\Delta_b$ is integral as a power series over $A_r$, where $\Delta_b=\prod_{c\ne b}(u_b-u_c)$. Here the inverse is taken in $A_r((X))$.

\begin{theorem}[Sharpness and cancellation]\label{thm:branches}
Observations recover the jet if and only if $n_b\ge1+t_b(r)$ for every $b$. If $H_*=0$ then $t_b(r)=d-1$. Otherwise
\[
 t_b(r)\le\omega+(r-1)D,\qquad
 t_b(r)=\omega+(r-1)D
 \ \Longleftrightarrow\
 p\nmid {M+r-2\choose r-1}.
\]
The older bound $\omega+(r-1)(q^{H_*}-1)$ is exact for $r=1$ or $H_*=0$. At $r>1$ and $H_*>0$ it is exact precisely when $H_*=1$ and $(M-1)+(r-1)$ has no carry in base $p$.
\end{theorem}

\begin{proof}
We first recall the interpolation argument in \cite{oldbranches}, including its necessity direction. Let $P_b(Y)=\prod_{c\ne b}(Y-u_c)$ and $P(Y)=\prod_c(Y-u_c)$. Every $u_c$ has zero constant term. The remainder $R$ on division of $H$ by $P$ in $Y$ has the same total-degree $<d$ jet, since $P\in(X,Y)^d$. Formal division is valid over the complete ring $A_r[[X]]$; the coefficients below the leading coefficient of $P$ belong to $(X)$. Over the Laurent ring, the differences $u_b-u_c$ are units, as also follows from the factorization below. Hence
\[
 R=\sum_b H(X,u_b)\,P_b(Y)/\Delta_b .
\]
If the observations vanish and $n_b\ge1+t_b$, each quotient $H(X,u_b)/\Delta_b$ lies in $XA_r[[X]]$, so $R\in(X,Y)^d$. Conversely $(X^{t_b}/\Delta_b)P_b(Y)$ vanishes on all other branches and equals $X^{t_b}$ at $b$. Its coefficient of $Y^{d-1}$ has nonzero constant term by minimality of $t_b$ (if $t_b=0$, the empty product case is immediate). Thus one fewer observed coefficient at $b$ can lose information even with exact data on every other branch.

Choose the Lubin--Tate coordinate $f(X)=[\pi](X)=\pi X+X^q$ \cite{lt}. An integral invertible change of parameter substitutes a series $X$ times a unit for $X$, and transforms differences by additional units. It preserves this threshold and the total-degree jet filtration. In the canonical coordinate,
\[
 [b](X)-[c](X)=[b-c](X)\cdot U(X)
              = f^{\circ h_c}(X)\cdot V(X)
 \quad (U,V\in\OO[[X]]^\times).
\]
The first unit comes from the formal group law $F(S,T)-S=T(1+\cdots)$; the second uses that $[(b-c)/\pi^{h_c}]$ has invertible linear coefficient. Write
\[
 g_j(X)=\pi+\big(f^{\circ(j-1)}(X)\big)^{q-1},\qquad
 e_j=(q-1)q^{j-1}.
\]
The identity $f^{\circ j}=f^{\circ(j-1)}g_j$ implies that $\Delta_b$ is associated to the distinguished monic polynomial
\begin{equation}
        P^{(b)}(X)=X^{d-1}\prod_{j=1}^{H_*}g_j(X)^{m_j},
        \qquad \deg P^{(b)}=d-1+\sum_j e_jm_j=\omega .
        \label{eq:denom}
\end{equation}
Every $g_j$ is Eisenstein of degree $e_j$, has constant $\pi$, and is congruent to $X^{e_j}$ modulo $\pi$. Put $g_j=X^{e_j}(1+\pi U_j)$, where $U_j\in X^{-1}\OO[X^{-1}]$ and the term at $X^{-e_j}$ has coefficient $1$. Over $A_r((X))$ the reciprocal of \eqref{eq:denom} is
\begin{equation}
 X^{-\omega}\prod_j\sum_{k=0}^{r-1}
       (-\pi)^k {m_j+k-1\choose k} U_j^k .
 \label{eq:inverse}
\end{equation}
Zero multiplicities can be omitted. A term of total index $k$ has extra pole order at most $kD$, and all terms with $k\ge r$ vanish. When $D>0$, the sole contribution to exponent $-\omega-(r-1)D$ chooses index $r-1$ at level $H_*$ and the extreme monomial in every copy of $U_{H_*}$. Its coefficient is $(-\pi)^{r-1}{M+r-2\choose r-1}$. It survives precisely when the binomial is a unit. Multiplication by an integral power-series unit preserves the least exponent even over $A_r$. Kummer's carry criterion gives the last assertion since $D=q^{H_*}-1$ exactly at $H_*=1$. With $H_*=0$, \eqref{eq:denom} is $X^{d-1}$.
\end{proof}

\paragraph{Exact second precision and frequency of equality.}
Let $\nu_j=\#\{c\ne b:h_c=j\}$. At $r=2$ the answer, including every cancellation, is particularly simple:
\[
 t_b(2)=\omega+\max\big(\{0\}\cup\{e_j:j\ge1,\ p\nmid\nu_j\}\big).
\]
Indeed, $f^{\circ j}\equiv X^{q^j}+\pi X^{q^{j-1}}\bmod\pi^2$. This follows by induction since $\pi q\in\pi^2\OO$. The reciprocal of the product of the iterates is therefore $X^{-\omega}(1-\pi\sum_{j\ge1}\nu_jX^{-e_j})$ modulo $\pi^2$. At arbitrary $r$ the upper endpoint in Theorem~\ref{thm:branches} is attained with a predictable frequency. If the digits of $M-1$ are $a_0,\ldots,a_{L-1}$ in base $p$, then the proportion of qualifying $r$ in every block of $p^L$ consecutive precisions is $\prod_{i=0}^{L-1}(p-a_i)/p^L$, by independent digit choices in $r-1$.

\paragraph{Two cancellation patterns.}
For $K=\QQ_3$, two directions that differ by $9$ have $\omega=9$ and $D=6$. Their common cutoff is $10+6(r-1)$ at every $r$; the old coefficient of $r-1$ was $8$. For $K=\QQ_2$, take the branch $0$ among $B=\{0,2,6\}$. The depth list is $(1,1)$ and the denominator is associated to $X^2(2+X)^2$. The reciprocal terms have coefficients $(-2)^k(k+1)$. Hence $t_0(r)=4+r-1$ for odd $r$ and $4+r-2$ for even $r$. At precisions one and two the same five coefficients suffice at this branch.

\paragraph{Relation to an established exact formula.}
For $K=\QQ_p$, $\pi=p$, and two directions with contact depth $h\ge1$, our formula gives
$t=p^h+(r-1)(p-1)p^{h-1}$. In the multiplicative coordinate the denominator is associated to $(1+X)^{p^h}-1$. The least exponent annihilating $X$ in the quotient is the nilpotency index of a cyclic group-ring augmentation generator. This special case also follows from the Weisman--Wilson results as presented by Clark and Schauz \cite[Theorem 4.7]{cs}. More than two branches give products of division factors with multiplicities and allow the cancellation in Theorem~\ref{thm:branches}.

\subsection{Optimal directions and the cost of increasing precision}
For any monic distinguished polynomial $Q$ write $\tau_r(Q)$ for the least $n$ such that $X^n=0$ in $A_r[X]/Q$. Division by $Q$ yields the same remainder in the polynomial ring and the formal power-series ring. In particular $t_b(r)=\tau_r(P^{(b)})$. This yields a terminating exact algorithm in the cancellation cases: start with $1$, multiply by $X$, and reduce modulo $Q$ and $\pi^r$ until the remainder is zero. For $Q=1$ take $\tau_r(1)=0$.

\begin{theorem}[All precisions and every residue ball]\label{thm:opt}
Suppose $d$ directions may be chosen in a prescribed coset of $\pi^a\OO$, $a\ge0$. The smallest possible common cutoff is
\begin{equation}
       1+\tau_r\!\left(X^{d-1}\prod_{j\ge1}g_j^{\,u_j}\right),
 \qquad
 u_j=\begin{cases}
    d-1&j\le a,\\[2pt]
    \left\lfloor(d-1)/q^{j-a}\right\rfloor&j>a .
 \end{cases} \label{eq:minimax}
\end{equation}
The same set of directions is optimal simultaneously at every precision.
\end{theorem}
\begin{proof}
The quantity $\tau_r$ increases weakly under divisibility. In a coset containing $d$ directions, successively follow a most populous child in the tree of residue balls. Its size after $h$ splits is at least $\lceil d/q^h\rceil$. A branch $b$ on this nested path therefore has $m_j(b)\ge u_j$ for all $j$; eventually its ball is a singleton. Equation \eqref{eq:denom} gives the lower bound. To attain it, take the base-$q$ expansions of $0,\ldots,d-1$ and replace digits by a fixed list of residue representatives used as $\pi$-adic digits. Scale by $\pi^a$ and translate to the chosen ball. After $h$ splits every residue class contains at most $\lceil d/q^h\rceil$ directions. Hence $m_j(b)\le u_j$ for every $b,j$. Divisibility gives the matching upper bound.
\end{proof}
For example, for eight directions over $\QQ_2$ the optimal unrestricted cutoff is $2r+11$. If all eight directions must be even, the optimum is $4r+21$. Indeed the profiles in \eqref{eq:minimax} are $(3,1)$ and $(7,3,1)$; in both the deepest level has multiplicity one.

\begin{proposition}[Growth with $r$]\label{prop:period}
For any fixed branch with $H_*>0$, the function $t_b(r)-Dr$ is bounded and eventually periodic; in particular $\lim_{r\to\infty}t_b(r)/r=D$.
\end{proposition}
\begin{proof}
Write $P=P^{(b)}$, $n=\deg P$, $\Lambda=\OO[X]/P$ and $x=X\bmod P$. The condition $t=t_b(r)$ implies $x^t\in\pi^r\Lambda$. Evaluation at a root of $g_{H_*}$ of valuation $1/D$ gives $t\ge Dr$. Theorem~\ref{thm:branches} gives $t\le Dr+\omega-D$.

On $\Lambda\otimes K$, multiplication by $S=x^D/\pi$ has integral characteristic polynomial: its eigenvalues are $\alpha^D/\pi$ for the roots $\alpha$ of $P$, with multiplicities. Each root is zero or has valuation $1/e_j\ge1/D$. The coefficients of the characteristic polynomial of $S$ therefore belong to $\OO$. By Cayley--Hamilton $\Lambda'=\sum_{i=0}^{n-1} S^i\Lambda$ is an $S$-stable lattice, also stable under $x$. It contains $\Lambda$ and $\pi^{n-1}\Lambda'\subseteq\Lambda$. For $0\le j\le\omega-D$, the test $x^{Dr+j}\in\pi^r\Lambda$ is equivalent to $x^j S^r\in\Lambda$. Membership depends only on the class of $S^r$ in the finite module $\Lambda'/\pi^{n-1}\Lambda'$. The orbit of $1$ under $S$ in this module is eventually periodic. So is the smallest $j$ passing the test.
\end{proof}

\paragraph{A quantitative consequence for integer-valued polynomials.}
Berger and Sprang use the case $r=1$ to multiply coefficient polynomials \cite[Theorem 4.6]{bs}. Their notation is $c_{f,j}(t)=[X^j]f([t](X))$, with $f$ in the kernel of the Lubin--Tate trace operator $\psi$. The following gives an explicit row cutoff at any residue precision. Set $d=i+\ell+1$, choose $B\subset\pi\OO$ of size $d$ by Theorem~\ref{thm:opt}, and define $h_b(X)=f(X)g([b](X))$. By \cite[Lemma 4.5]{bs}, $h_b\in\ker\psi$ if $f\in\ker\psi$ and $g\in\OO[[X]]$. Define constants in $A_r$ by
\[
       \mu_{bj}=[X^{i-j}Y^\ell]\,P_b(Y)/\Delta_b .
\]
For every $t\in\OO$ the following congruence of integer-valued functions holds:
\begin{equation}
 c_{f,i}(t)c_{g,\ell}(t)\equiv
 \sum_{b\in B}\sum_{j=0}^{t_b(r)}\mu_{bj}\,c_{h_b,j}(t)\pmod{\pi^r}.
 \label{eq:products}
\end{equation}
Apply interpolation to $f([t](X))g([t](Y))$. The branch values are $h_b([t](X))$ since the endomorphisms commute. Moreover $[Y^\ell]P_b\in X^{d-1-\ell}A_r[[X]]$, which gives the displayed cutoff. Equation \eqref{eq:products} supplies explicit weights and an optimal common precision for the full jet used by this procedure. A single prescribed product can require fewer coefficients.

\section{The largest spanning-tree valuation in a rank-two graph tower}
\label{sec:graph-extreme}

How divisible by $p$ can the spanning-tree count of a connected abelian Cayley graph be, if the number of steps is fixed? Question 2 in the earlier graph-tower report asks for the largest coefficient $A$ when the quadratic moment vanishes identically. We determine the exact maximum for four steps at every prime $p>8$. For every even number of steps we give a sharp upper bound and characterize the primes where equality is possible. The earlier report treats the case of a nonzero quadratic moment \cite{oldgraphs}. DuBose and Valli\`eres prove the eventual-polynomial theorem for these towers; Dion, Lei, Ray and Valli\`eres study related extremal orders in one-variable towers \cite{dv,dlrv}. The companion \emph{Effective Graph Towers} report addresses the starting level of the polynomial \cite{effective}.

\subsection{The cyclotomic calculation}
Let $p>2m$ be prime. The $m$ integral step vectors $s_i$ span $\mathbb F_p^2$ after reduction. Repetitions are allowed. At level $n\ge1$ the vertices are $(\mathbb Z/p^n)^2$ and each pair $(v,i)$ supplies an undirected edge from $v$ to $v+s_i$. Write $T_n$ for the $p$-adic valuation of the spanning-tree count, normalized by $v_p(p)=1$. Define binary forms and a projective depth by
\[
 M_j(u)=\sum_{i=1}^m\langle u,s_i\rangle^{2j},
 \qquad r(L)=\min\{j\ge1: M_j(u)\ne0\},\qquad L=[u]\in\mathbb P^1(\mathbb F_p).
\]
Here $M_1$ is the quadratic moment, and the forms have coefficients in $\mathbb F_p$. Newton's identities give $r(L)\le m$. If the first $m$ power sums of the squared projections vanished, every projection would be zero, contradicting spanning. The prior packet proves
\begin{equation}
 A=2\sum_L r(L),\qquad T_n=A(1+p+\cdots+p^{n-1})-2n.       \label{eq:graph-coeff}
\end{equation}
We recall its derivation to identify where torsion specialization enters. On a character of conductor $p^h$, with primitive covector $u$ modulo $p^h$, the Laplacian eigenvalue is $\sum_i(2-Z^{a_i}-Z^{-a_i})$, where $a_i=\langle u,s_i\rangle$ and $Z=\zeta_{p^h}$. Write $Z=1+T$. In characteristic $p$ the truncated substitution $T=\exp(z)-1$ is invertible through degree $2m<p$. The coefficient of $z^{2j}$ is $-2M_j(u)/(2j)!$, with odd coefficients zero. The first unit coefficient in the original integral $T$-series is therefore at $k=2r([\bar u])$. At $T=\zeta_{p^h}-1$, the valuation is $k/e$, $e=\varphi(p^h)$. Indeed $k\le p-1\le e$, the constant coefficient is zero, and every earlier nonzero coefficient is divisible by $p$. Thus the term of order $k$ is the unique term of least valuation. This uses the torsion radius in \src{New/Congruences/Selmer/RootsOfUnity.lean\#L58}{\texttt{ellVal\_sub\_one\_eq\_rootRadius}}.

For each residual direction $L$, normalize one unit coordinate to one. There are $p^{h-1}$ projective lifts to conductor $p^h$, each with $e$ conjugate characters. This direction contributes $2r(L)p^{h-1}$ to the valuation of the product. The matrix-tree theorem divides the nonzero-eigenvalue product by the vertex count $p^{2n}$, proving \eqref{eq:graph-coeff}.

\subsection{An interpolation bound and an attaining configuration}
Put $s=\lfloor m/2\rfloor$. For every spanning list,
\begin{equation}
 A\le 2s(p+1)+4s(m-s).                                      \label{eq:graph-bound}
\end{equation}
For even $m=2s$, equality is attainable exactly at primes $2m\mid p-1$. The maximizing construction has $M_1=0$ whenever $m\ge4$.

\emph{Proof.} Let $t$ be the first $j$ with $M_j$ a nonzero polynomial. Its degree is $2t\le 2m<p$. We prove $t\le s$. Discard residue-zero steps and let $m'$ be the remaining count. If a line through the origin contains more than half of these steps, take a nonzero linear functional annihilating the line. Its list has at most $\lfloor m'/2\rfloor$ nonzero projections and at least one, by spanning. Newton's identities applied to the shorter list prove the bound.

Otherwise, arrange the $2m'$ signed occurrences $\pm s_i$ in contiguous blocks according to the line through the origin that they occupy. Every block has size at most $m'$. Pair occurrence $j$ with occurrence $j+m'$ for $1\le j\le m'$. Each pair has two distinct radial directions. Its affine line avoids the origin. Multiplying the equations of these $m'$ affine lines produces a polynomial $H$ of degree $m'$ that vanishes on every signed occurrence and satisfies $H(0)\ne0$.

If $t>\lfloor m'/2\rfloor$, then $m'\le2t-1$. For every polynomial $P$ of degree at most $2t-1$, we have
\[
 \sum_{i:s_i\ne0}\big(P(s_i)+P(-s_i)\big)=2m'P(0).
\]
Odd homogeneous degrees cancel. For degree $2j<2t$, $j\ge1$, the vanishing coefficients of $M_j$ say that each mixed monomial of that degree sums to zero, since all binomial coefficients $\binom{2j}{k}$ are units. Applying the identity to $H$ contradicts $p>2m$. This proves $t\le s$ in both cases.

A nonzero binary form of degree $2t$ has at most $2t$ projective roots. The other directions have depth $t$, and each exceptional direction has depth at most $m$. Hence $A/2\le t(p+1)+2t(m-t)$. The right side increases strictly for integers $1\le t\le s$, which proves \eqref{eq:graph-bound}.

For even $m=2s$, equality requires $2s$ exceptional directions of depth $m$. In any direction of depth $m$, Newton's identities give $\prod_i (X-\langle u,s_i\rangle^2)=X^m-c$, $c\ne0$. The roots are distinct nonzero squares. Dividing by one root shows that all $m$th roots of unity belong to the square subgroup, and hence $2m\mid p-1$.

Conversely, assume $4s\mid p-1$ and choose $g$ of order $2s$. Use the $s$ steps $(g^j,0)$ and the $s$ steps $(0,g^j)$, $0\le j<s$, lifting residues to integers. Every moment of index $k$ with $s\nmid k$ vanishes identically by a geometric sum. The moment $M_s$ is $s(a^{2s}+b^{2s})$ in direction $[a:b]$. Exactly $2s$ rational directions solve $a^{2s}=-b^{2s}$ because $-1$ is a $2s$th power. In each, $M_{2s}=2s b^{4s}\ne0$, so their depth is $2s=m$. All others have depth $s$. Substitution proves equality.\hfill$\square$

For $p=13$ take three steps on each axis with coordinate sizes $1,4,3$. The eight ordinary directions have depth three, and the six exceptional directions have depth six. Thus $A=120$, $T_1=118$, and $T_2=1676$. For this same prime and step count the earlier nonvanishing-quadratic maximum was $48$. The included program evaluates the $168\times168$ reduced integral Laplacian determinant directly and recovers $T_1=118$.

\subsection{Four steps at every prime}
For $m=4$ and every prime $p>8$, the unrestricted maximum, also the maximum with $M_1=0$, is
\begin{equation}
 A_{\max}=4(p+1)+
 \begin{cases}
  16,&p\equiv1\pmod8,\\
  0,&p\in\{11,23,29,47,53,83,149\},\\
  8,&\text{otherwise}.
 \end{cases}                                                    \label{eq:graph-four}
\end{equation}

\emph{Proof.} The primes $p\equiv1\pmod8$ are covered by the extremal construction. If $M_1$ is nonzero, the at-most-two exceptional directions give $A\le2(p+1)+12<4(p+1)$. It remains to optimize over lists with $M_1=0$ at the other primes. Relabel two independent steps and make them $e_1,e_2$ by a change of basis. The two rows of the remaining $2\times2$ block must be perpendicular, both with sum of squared coordinates $-1$. The first row $(a,b)$ determines the second up to sign. Changing the sign of the second basis coordinate, followed by a sign change of the second step, if needed, puts the block in the form
\[
 B=\begin{pmatrix}a&b\\b&-a\end{pmatrix},\quad a^2+b^2=-1,
 \qquad J=\begin{pmatrix}0&-1\\1&0\end{pmatrix}.
\]
Such $a,b$ exist for every odd prime because the squares and $-1$ minus the squares each have $(p+1)/2$ elements. The bound on $t$ gives $t=2$.

Both $B$ and $J$ permute the signed steps $\{\pm e_i,\pm Be_i:i=1,2\}$, with multiplicities. They satisfy $J^2=B^2=(JB)^2=-I$ and $JB=-BJ$. Their projective transformations on covectors generate a group of order four with three involutions. Indeed $B$ is symmetric and nonscalar, $J$ is skew-symmetric and invertible, and $JB$ is projectively different from both. Every $M_j$ is invariant under the corresponding transpose matrices. A direction fixed by any nonidentity element has eigenvalue $\lambda$ with $\lambda^2=-1$ for that representative. Then $M_3(u)=\lambda^6M_3(u)=-M_3(u)$. If such a direction were a root of the quartic $M_2$, it would have depth four, which requires $8\mid p-1$. Thus all rational roots of $M_2$ are in orbits of four. Since $M_2$ is a nonzero quartic, it has either zero or four rational roots, and each has depth three. This contributes precisely zero or eight above $4(p+1)$.

Existence of at least one such root, allowing the step list to vary, is equivalent to a rational point on the projective curve
\[
 C:\quad x_1^2+x_2^2+x_3^2+x_4^2
       =x_1^4+x_2^4+x_3^4+x_4^4=0 .
\]
The four projections in an exceptional direction give such a point. Conversely, take a nonzero representative $x$ of a point of $C(\mathbb F_p)$ and put $D=\operatorname{diag}(J,J)$. The vectors $x,Dx$ are perpendicular and isotropic in $\mathbb F_p^4$. If $Dx=\lambda x$, invariance under $D$ gives $\sum x_i^6=0$ as well, since $\lambda^2=-1$. The Newton argument for the four squares $x_i^2$ then forces $8\mid p-1$. For the primes under consideration the rows $x,Dx$ are therefore independent. Their four columns supply spanning steps with identically zero quadratic moment, and the first coordinate functional gives the required quartic root.

Over an algebraic closure $C$ is a smooth complete intersection. Dependence of its two gradients would force all nonzero $x_i^2$ to coincide with one common nonzero $c$. The quadratic equation would give $kc=0$ for some $1\le k\le4<p$. The quartic does not vanish identically on the smooth quadric, as the point $(1,\sqrt{-1},0,0)$ shows. On this quadric, isomorphic over the algebraic closure to $\mathbb P^1\times\mathbb P^1$, the curve has bidegree $(4,4)$. It is connected. If two nonempty effective divisors of classes $(a,b)$, $(4-a,4-b)$ partitioned it disjointly, their intersection number $a(4-b)+b(4-a)$ would be positive. Smoothness then gives geometric irreducibility. Adjunction gives $g=(4-1)(4-1)=9$. The classical Hasse--Weil theorem \cite{weil} yields
\[
 \#C(\mathbb F_p)\ge p+1-18\sqrt p>0\qquad (p>322).
\]
This use of the finite-field theorem is independent of Dirichlet GRH.

For the 49 remaining primes $8<p\le322$, $p\not\equiv1\pmod8$, normalize a nonzero coordinate to one using permutation and scaling. Loop through $y,z\in\mathbb F_p$. Set $c=-1-y^2-z^2$ and test whether $c$ is a square and $1+y^4+z^4+c^2=0$. This exhausts $C(\mathbb F_p)$ for the existence test. The seven failures are exactly those in \eqref{eq:graph-four}. The accompanying code records a witness for every successful prime and independently enumerates all plane normal forms at eleven small primes.\hfill$\square$

\paragraph{Simple graphs.} The four-step maximum is also attained among simple eight-regular Cayley graphs. The row plane $W\subset\FF_p^4$ of a spanning list with $M_1=0$ is totally isotropic for the dot product. A zero step or equality of two steps up to sign would put $W$ in the orthogonal complement of $e_i$ or $e_i\pm e_j$. Each complement is a nondegenerate three-space, in which a totally isotropic plane would violate $\dim W+\dim W^\perp=3$. Distinct nonzero steps up to sign modulo $p$ give a simple graph at every level.

The interpolation bound leaves open the exact maxima for odd $m\ge5$ and for even $m\ge6$ with $2m\nmid p-1$. Simultaneous projective roots of several moment forms determine which exceptional depths can coexist on a single plane.

\section{How far does the one-sided zero test transfer?}
The local applications above require no hypothesis on zeros. The global strategy faces a different question. Which prime sums can be bounded by the source's determinant construction? A test case from the incoming Frobenius report \cite{oldcubic} is the splitting field $L$ of $X^3-X-1$. The cubic is irreducible modulo $2$ and its discriminant is $-23$, so its group is $S_3$ and its quadratic subfield is $K=\QQ(\sqrt{-23})$. For an unramified prime $p$, write $I(p)=1$ if the cubic remains irreducible modulo $p$, and $0$ otherwise.

\paragraph{The missing inequality.}
With Dirichlet GRH in hand, the earlier criterion for RH for $\zeta_L$ becomes the following one-sided estimate, for each $\epsilon>0$:
\begin{equation}
        \sum_{p\le x,\ p\ne23} I(p)\log p-\frac{x}{3}
           \ \le\ C_\epsilon x^{1/2+\epsilon}.                    \label{eq:cubic}
\end{equation}
In $S_3$ the three-cycle indicator is $\mathbf1_{A_3}-\mathbf1_{\{1\}}$. Add higher unramified prime powers with this weight on Frobenius powers. The resulting $\Psi_I(x)$ differs from the prime sum by $O(\sqrt{x}\log^2(2x))$. For $\Re s>1$ the Mellin transform of $\Psi_I(x)-x/3$ against $x^{-s-1}dx$ on $[1,\infty)$ is
\[
 \frac1s\left(-\frac12\frac{\zeta'_K}{\zeta_K}
                 +\frac16\frac{\zeta'_L}{\zeta_L}+U(s)\right)
                 -\frac{1}{3(s-1)},
\]
where the omitted ramified factors contribute $U$ holomorphic for $\Re s>0$. This follows by comparing the Euler factors over each unramified prime. The pole at $1$ cancels. Dirichlet GRH gives zero-freeness of $\zeta_K$ for $\Re s>1/2$ away from its pole, so each zero $\rho$ of $\zeta_L$ to the right of the line produces a residue $\ord_\rho(\zeta_L)/(6\rho)$.

If \eqref{eq:cubic} holds, choose $\theta=1/2+\epsilon<1$ and a constant $C$ making $Cx^\theta-\Psi_I(x)+x/3$ nonnegative. Its Mellin transform has a meromorphic continuation to $\Re s>\theta$ whose possible real poles there have strictly negative residue. Landau's theorem forces a real singularity at a finite abscissa of convergence, whereas such a negative residue would make the transform negative immediately to the right. Thus the abscissa is at most $\theta$ and the integral is holomorphic for $\Re s>\theta$. Taking $\epsilon$ arbitrarily small and applying the functional equation proves the desired RH. Conversely RH for $L$ and $K$ gives \eqref{eq:cubic} by their prime ideal theorems \cite{gm}. This recovers the specific prior reduction and locates its unproved input.

\paragraph{The exact error of a fixed congruence approximation.}
Suppose one tries to express the weight $I(p)$ by a real function $w$ of $p\bmod Q$, for a fixed modulus $Q\ge1$. The best possible mean square error is
\begin{equation}
 \min_w\ \lim_{x\to\infty}\frac{1}{\#\{p\le x\}}
 \sum_{\substack{p\le x\\p\nmid23Q}}\!\big(I(p)-w(p\bmod Q)\big)^2
   =
 \begin{cases} 2/9,&23\nmid Q,\\ 1/9,&23\mid Q. \end{cases}       \label{eq:approx}
\end{equation}
Indeed, $L\cap\QQ(\zeta_Q)$ is abelian over $\QQ$, hence is either $\QQ$ or $K$. The quadratic field has conductor $23$, so the second case occurs exactly for $23\mid Q$. Chebotarev's theorem for the compositum gives uniform Frobenius distributions within the compatible fibers for conjugacy-invariant tests \cite{cheb}. If the intersection is trivial, the conditional probability of a three-cycle in every reduced class is $1/3$. Otherwise half the classes force odd parity, with probability zero of a three-cycle, and half force even parity, with probability $2/3$. In each class the minimizer is its conditional probability and the error is the Bernoulli variance. This gives \eqref{eq:approx}. Even the best binary congruence rule, with the modulus freely chosen, misclassifies a proportion $1/6$ of primes. It predicts irreducibility on even classes and reducibility on odd classes.

This calculation excludes an exact replacement of the prime weight by any fixed collection of Dirichlet characters. The error in \eqref{eq:cubic} is a signed sum and could still exhibit cancellation. Establishing that cancellation by a determinant estimate for the three-cycle weight would supply the missing analytic input.

\section{Verification and further questions}
In the four-step classification a finite enumeration is part of the proof. Two implementations reproduce the seven-prime list. One searches for a normalized witness and one counts the entire curve by convolving two-coordinate moment tables. The latter checks all 62 primes from $11$ through $317$. Separately, reduced Laplacians at $p=11,13$ give four-step valuations $46,62$, and an integer cofactor gives the six-step valuation $118$ at $p=13$.

For branch precision, a monic-remainder computation checks 1,400 combinations of contact profiles and precisions. A separate Laurent computation compares 640 cases with literal branches $(1+X)^b-1$. A third check constructs the observation matrices over $\ZZ/p^r$ and compares image sizes before and after adjoining the desired Taylor coefficients. It accepts eight sufficient cutoff vectors and detects loss of information at every one of their 25 one-coordinate reductions. A further 6,128 finite residue-ball designs test the minimax formula. These finite checks illustrate the general interpolation proofs. The new theorems are proved here on paper and remain to be formalized in Lean. Publication priority remains unverified.

\paragraph{Further research questions.}
For the branch problem, determine the least eventual period of $t_b(r)-Dr$ directly from the contact multiplicities. For graph towers, which root configurations of the quartic, sextic and higher directional moments can occur on a single plane for $m=5$? For the cubic splitting field, can a geometric construction produce \eqref{eq:cubic}? Formula \eqref{eq:approx} quantifies the information that would have to come from outside any fixed congruence-class description.

\begin{thebibliography}{99}\small
\bibitem{repo} L. Sernau and collaborators, \emph{GRH for Dirichlet $L$-functions, in Lean 4}. Shorter rewriting of the Alp\"oge--Claude proof. \href{https://github.com/lsernau/grh-lean/tree/72fead102f663d5c7a407964a13dc0cf9d977207}{Commit 72fead102f66}, 25 November 2026.
\bibitem{sutherland} A. V. Sutherland, \emph{Verification files for the shorter Lean proof of RH and of GRH for Dirichlet $L$-functions}, 26 November 2026. \href{https://math.mit.edu/~drew/shortgrh/}{Test logs and challenge statements}.
\bibitem{oldbranches} ProveIt incoming, \emph{Optimal signings of square cycles and sharp Lubin--Tate jet recovery}. Integral-precision proposition and further question 3. \href{https://github.com/VladimirReshetnikov/ProveIt/blob/1677488425e45c3df14149c804014133798136a5/docs/incoming/grh_transfers_cycle_signings_and_lubin_tate.zip}{Pinned manuscript archive}.
\bibitem{oldgraphs} ProveIt incoming, graph towers and finite rings report, \texttt{moments.tex} and \texttt{endmatter.tex}, question 2. \href{https://github.com/VladimirReshetnikov/ProveIt/blob/1677488425e45c3df14149c804014133798136a5/docs/incoming/grh_method_transfers_graph_towers_and_finite_rings.zip}{Pinned manuscript archive}.
\bibitem{effective} ProveIt incoming, \emph{When does the spanning-tree formula start? Effective Iwasawa polynomials from cyclotomic valuations}. \href{https://github.com/VladimirReshetnikov/ProveIt/blob/1677488425e45c3df14149c804014133798136a5/docs/incoming/Effective_Graph_Towers.zip}{Pinned manuscript archive}.
\bibitem{oldcubic} ProveIt incoming, McNew density and GRH transfers report, \texttt{frobenius.tex}. \href{https://github.com/VladimirReshetnikov/ProveIt/blob/1677488425e45c3df14149c804014133798136a5/docs/incoming/mcnew_density_and_grh_transfers.zip}{Pinned manuscript archive}.
\bibitem{bs} L. Berger and J. Sprang, \emph{Integer-valued polynomials and $p$-adic Fourier theory}, \S4. \href{https://arxiv.org/html/2502.18053v2}{arXiv:2502.18053v2}.
\bibitem{lt} J. Lubin and J. Tate, \emph{Formal complex multiplication in local fields}, Ann. of Math. \textbf{81} (1965), 380--387, Theorems 1--2. \href{https://doi.org/10.2307/1970622}{Publisher}.
\bibitem{cs} P. L. Clark and U. Schauz, \emph{Functional degrees and arithmetic applications, I}, Theorem 4.7 and its antecedents due to Weisman and Wilson. \href{https://arxiv.org/abs/2201.02763}{arXiv:2201.02763}.
\bibitem{dv} S. DuBose and D. Valli\`eres, \emph{On $\ZZ_\ell^d$-towers of graphs}, Algebr. Comb. \textbf{6} (2023), 1331--1346. \href{https://doi.org/10.5802/alco.304}{Publisher}.
\bibitem{dlrv} C. Dion, A. Lei, A. Ray and D. Valli\`eres, \emph{On the distribution of Iwasawa invariants associated to multigraphs}, Nagoya Math. J. \textbf{253} (2024), 48--90. \href{https://doi.org/10.1017/nmj.2023.18}{Publisher}.
\bibitem{weil} A. Weil, \emph{Sur les courbes alg\'ebriques et les vari\'et\'es qui s'en d\'eduisent}, Hermann, 1948. \href{https://denisevellachemla.eu/Weil-courbes-varietes.pdf}{Monograph}.
\bibitem{gm} L. Greni\'e and G. Molteni, \emph{Explicit versions of the prime ideal theorem for Dedekind zeta functions under GRH}, Math. Comp. \textbf{85} (2016), 889--906. \href{https://air.unimi.it/handle/2434/448115}{Article record}.
\bibitem{cheb} J. C. Lagarias, H. L. Montgomery and A. M. Odlyzko, \emph{A bound for the least prime ideal in the Chebotarev density theorem}, Invent. Math. \textbf{54} (1979), 271--296. \href{https://doi.org/10.1007/BF01390234}{Publisher}.
\end{thebibliography}

\end{document}
